\section{Proof of \texorpdfstring{Theorem~\ref{thm:erw-bounds}}{Theorem 2.6}}\label{app:erw}

Throughout, $N\ge2$ is even and we condition on $X_1=+1$. By
Proposition~\ref{prop:first-mutation}(a) we need to bound $V_{t+1}(1)$, and we
compare $V$ with explicit hitting probabilities of P\'olya urns. Let
$\mathcal L_sF(b)=(1-r_s(b))F(b)+r_s(b)F(b+1)$, so that $V_s=\mathcal L_sV_{s+1}$ and $V_N(b)=\mathbf 1\{b=h\}$. Both $M_t$
and $t-M_t\ge1$ are nondecreasing. So from $b\ge1$ the chain stays at states
$\ge1$, and for $b\ge1$ we have $V_s(b)=0$ unless $b$ is feasible, that is,
$b\in\mathcal S_s=\{b\in\Z:1\le b\le h,\ 1\le s-b\le h\}$. Note that
$\mathcal S_N=\{h\}$.

For $c\ge0$ the P\'olya$(c)$ chain moves at time $s$ from $b$ to $b+1$ with
probability $r^c_s(b)=\frac{b+c}{s+2c}$. It is P\'olya's urn with weights
$b+c$ and $s-b+c$ for the 2 colors, that is, with $c$ extra balls of each
color. For $b+c>0$ and $s-b+c>0$ let
$\Phi^c_s(b)=\binom{N-s}{h-b}\mathrm B(h+c,h+c)/\mathrm B(b+c,s-b+c)$, where
$\mathrm B$ is the beta function and the binomial coefficient is 0 unless
$0\le h-b\le N-s$. Theorem~\ref{thm:erw-bounds}(a) will follow from
Lemma~\ref{lem:A-polya}(d), and (b) from Proposition~\ref{prop:A-super}.

\begin{lemma}\label{lem:A-polya}
Let $c\ge0$, $1\le s\le N$, $b+c>0$ and $s-b+c>0$.
\begin{enumerate}[(a)]
\item $\Phi^c_s(b)$ is the probability that the P\'olya$(c)$ chain at $b$ at
time $s$ is at $h$ at time $N$. So $\Phi^c_N(b)=\mathbf 1\{b=h\}$, and
$\Phi^c_s(b)=r^c_s(b)\Phi^c_{s+1}(b+1)+(1-r^c_s(b))\Phi^c_{s+1}(b)$ for
$s<N$.
\item If $s<N$ and $\Phi^c_s(b)>0$, then
$r^c_s(b)\Phi^c_{s+1}(b+1)/\Phi^c_s(b)=\bar r_s(b):=\frac{h-b}{N-s}$.
\item If $0<\varepsilon<\frac12$ and $\beta=\frac{\varepsilon}{1-2\varepsilon}$,
then $r_s(b)=r^{\beta s}_s(b)$ for $0\le b\le s$.
\item If $0<\varepsilon<\frac12$ and $0\le c\le\beta s$, then
$V_s(b)\ge\Phi^c_s(b)$ for $1\le b\le s-1$. In particular
$V_{t+1}(1)\ge\Phi^0_{t+1}(1)=\pi_t$ for $1\le t\le h$.
\end{enumerate}
\end{lemma}

\begin{proof}
(a) Let $x=b+c$, $y=s-b+c$ and $(x)_k=x(x+1)\cdots(x+k-1)$. As in the proof
of Lemma~\ref{lem:urn}, each order of the $N-s$ draws with $h-b$ draws of
the first color has probability
$(x)_{h-b}(y)_{N-s-h+b}/(x+y)_{N-s}=\mathrm B(h+c,h+c)/\mathrm B(x,y)$, and
there are $\binom{N-s}{h-b}$ such orders. The recursion is the Markov
property.

(b) Use $\binom{N-s-1}{h-b-1}/\binom{N-s}{h-b}=\frac{h-b}{N-s}$ and
$\mathrm B(x,y)/\mathrm B(x+1,y)=\frac{x+y}x=1/r^c_s(b)$.

(c) Since $1+2\beta=1/(1-2\varepsilon)$, we have
$\frac{b+\beta s}{s+2\beta s}=(1-2\varepsilon)(\frac bs+\beta)=r_s(b)$.

(d) Let $F_u=\Phi^c_u$ on $\{1,\dots,u-1\}$ for $s\le u\le N$. Then $F_u$
vanishes off $\mathcal S_u$ and $F_N(b)=\mathbf 1\{b=h\}$. By (a) and (c),
$\mathcal L_uF_{u+1}(b)-F_u(b)=(r_u(b)-r^c_u(b))(F_{u+1}(b+1)-F_{u+1}(b))$.
Since $r^c_u(b)=\frac12-(\frac u2-b)/(u+2c)$, this is the product of
$\frac u2-b$, $\frac1{u+2c}-\frac1{u(1+2\beta)}\ge0$ (as $c\le\beta s\le\beta
u$) and $F_{u+1}(b+1)-F_{u+1}(b)$. Let
$b\in\mathcal S_u$. If $b=h$ then $b\ge u/2$ and $F_{u+1}(b+1)=0$. If
$u-b=h$ then $b\le u/2$ and $F_{u+1}(b)=0$. Otherwise
$F_{u+1}(b+1)/F_{u+1}(b)=\frac{(h-b)(u-b+c)}{(h-u+b)(b+c)}$, which is
$\ge1$ exactly when $b\le u/2$ (write $b=\frac u2-d$; the difference of
the 2 sides of the inequality is $2d(h+c)$). So in all 3 cases the first and
last factors have the same sign, and $\mathcal L_uF_{u+1}\ge F_u$ on
$\mathcal S_u$, and Lemma~\ref{lem:A-compare}(b) gives $V_s\ge F_s$.
Finally $\Phi^0_{t+1}(1)=\beta_{t+1}(h)=\pi_t$ by (a) and
Lemma~\ref{lem:urn}.
\end{proof}

\begin{lemma}\label{lem:A-compare}
Let $1\le s_0\le N$, and for $s_0\le s\le N$ let $F_s\ge0$ be a function on
$\{1,\dots,s-1\}$.
\begin{enumerate}[(a)]
\item If $F_N(h)\ge1$ and $\mathcal L_sF_{s+1}\le K_sF_s$ on $\mathcal S_s$
for $s_0\le s<N$, with constants $K_s\ge0$, then
$V_s\le\bigl(\prod_{u=s}^{N-1}K_u\bigr)F_s$ on $\mathcal S_s$.
\item If $F_N(b)\le\mathbf 1\{b=h\}$, each $F_s$ vanishes off $\mathcal S_s$,
and $\mathcal L_sF_{s+1}\ge F_s$ on $\mathcal S_s$ for $s_0\le s<N$, then
$V_s\ge F_s$.
\end{enumerate}
\end{lemma}

\begin{proof}
We use backward induction on $s$, from $V_N(b)=\mathbf 1\{b=h\}$. In (a),
$V_{s+1}$ vanishes at infeasible states and $F_{s+1}\ge0$, so on
$\mathcal S_s$ we get
$V_s=\mathcal L_sV_{s+1}\le\prod_{u>s}K_u\,\mathcal L_sF_{s+1}\le\prod_{u\ge s}K_u\,F_s$.
In (b) we get $V_s\ge\mathcal L_sF_{s+1}\ge F_s$ on $\mathcal S_s$, and
$F_s=0\le V_s$ off $\mathcal S_s$.
\end{proof}

For the upper bound we compare $M$ with an urn that has more extra balls,
and we control the cost of each step through the derivative of $\Phi^\tau_s$
in $\tau$. Let $\psi$ be the digamma function.

\begin{lemma}\label{lem:A-cderiv}
Let $2\le s\le N$, $b\in\mathcal S_s$ and $\tau\ge0$. Put $X=b+\tau$,
$Y=s-b+\tau$, $\Sigma=X+Y$ and $z=(s-2b)^2/\Sigma^2$. Then
\[
 \partial_\tau\log\Phi^\tau_s(b)
 <-\log(1-z)+\frac{1+z}{(1-z)\Sigma}+\frac{2(1+z)}{3\Sigma^2(1-z)^2}.
\]
\end{lemma}

\begin{proof}
Since $\log\mathrm B(x,y)=\log\Gamma(x)+\log\Gamma(y)-\log\Gamma(x+y)$, the
derivative is $2\psi(\Sigma)-\psi(X)-\psi(Y)-2\psi(N+2\tau)+2\psi(h+\tau)$.
For $x>0$, \cite[5.9.15]{dlmf} writes $\log x-\frac1{2x}-\psi(x)$ as
$2\int_0^\infty\frac{t\,dt}{(t^2+x^2)(e^{2\pi t}-1)}$, which lies in
$(0,\frac1{12x^2})$ since $\int_0^\infty\frac{t\,dt}{e^{2\pi t}-1}=\frac1{24}$.
We use this upper bound for $\psi(\Sigma)$ and lower bound for $\psi(X)$ and
$\psi(Y)$. The duplication formula \cite[5.5.5]{dlmf} gives
$2\psi(2y)=\psi(y)+\psi(y+\frac12)+2\log2$, and $\psi$ is increasing, so
$2\psi(N+2\tau)-2\psi(h+\tau)>2\log2$. Hence the derivative is less than
$\log\frac{\Sigma^2}{4XY}+(\frac1{2X}+\frac1{2Y}-\frac1\Sigma)
+\frac1{12}(\frac1{X^2}+\frac1{Y^2})$. Now $4XY=\Sigma^2(1-z)$ and
$X^2+Y^2=\Sigma^2(1+z)/2$, so the 3 terms are $-\log(1-z)$,
$\frac{X^2+Y^2}{2XY\Sigma}=\frac{1+z}{(1-z)\Sigma}$ and
$\frac{X^2+Y^2}{12X^2Y^2}=\frac{2(1+z)}{3\Sigma^2(1-z)^2}$.
\end{proof}

\begin{proposition}\label{prop:A-super}
Let $0<\varepsilon<\frac14$, $\kappa=\frac{2\varepsilon}{1-4\varepsilon}$,
$\kappa_s=\kappa(s+10)$ and $\Sigma_s=s+2\kappa_s$. For $1\le s\le N-1$ and
$b\in\mathcal S_s$ let
$\Delta_s(b)=\mathcal L_s\Phi^{\kappa_{s+1}}_{s+1}(b)\big/\Phi^{\kappa_s}_s(b)$.
\begin{enumerate}[(a)]
\item Let $c=\kappa_s$, $c'=\kappa_{s+1}$, $d=\frac s2-b$, $u=\frac{d}{s+2c'}$,
$v=\frac d{N-s}$ and $\rho=\frac{s+2c'}{s(1+2\beta)}$. Then $|u|<\frac12$
and
\[
 \Delta_s(b)=\mathcal J\,\frac{\Phi^{c'}_s(b)}{\Phi^{c}_s(b)},\qquad
 \mathcal J=1-(\rho-1)\frac{u(u+v)}{\frac14-u^2}.
\]
\item $\log \Delta_s(b)\le\frac{\kappa}{\Sigma_s}+\frac{2\kappa}{3\Sigma_s^2}$.
\item For $1\le t\le h$,
$V_{t+1}(1)\le\pi_t\exp\bigl(\kappa\bigl[H_{N-1}-H_t+\frac2{3t}\bigr]
+\kappa_{t+1}\bigl[\log\frac{(t+1)^2}{4t}+\frac23\bigr]\bigr)$.
\item $C_N\le U_N:=\varepsilon c_1e^{\kappa H_{N-1}}\sum_{t=1}^hw_t
(1-\varepsilon)^{t-1}e^{\kappa\mathcal G(t)}$, where
$\mathcal G(t)=(t+11)\bigl(\log\frac{(t+1)^2}{4t}+\frac23\bigr)-H_t+\frac2{3t}$.
\end{enumerate}
\end{proposition}

The slope $\kappa$ solves $2(\kappa-\beta)/(1+2\beta)=\kappa$, so that in (b)
the gain in $\mathcal J$ cancels the main loss from raising the extra balls
from $c$ to $c'$.

\begin{proof}
(a) Write $r'=r^{c'}_s(b)$, $\bar r=\bar r_s(b)$ and $r=r_s(b)$. By
Lemma~\ref{lem:A-polya}(b) with $c'$, and the recursion in
Lemma~\ref{lem:A-polya}(a), we have
$r'\Phi^{c'}_{s+1}(b+1)=\bar r\,\Phi^{c'}_s(b)$ and
$(1-r')\Phi^{c'}_{s+1}(b)=(1-\bar r)\Phi^{c'}_s(b)$. So the numerator of
$\Delta_s(b)$ divided by $\Phi^{c'}_s(b)$ is
$\frac{(1-r)(1-\bar r)}{1-r'}+\frac{r\bar r}{r'}$. Here $r'=\frac12-u$,
$\bar r=\frac12+v$ and $r=\frac12-\rho u$, since
$\frac12-r=(1-2\varepsilon)\frac ds=\rho u$. Over the denominator
$r'(1-r')=\frac14-u^2$ the numerator is $f(u,v)+f(-u,-v)$ with
$f(u,v)=(\frac12+\rho u)(\frac12-v)(\frac12-u)$, and this equals
$\frac14-\rho u^2-(\rho-1)uv$. Also $|d|\le\frac s2-1$, so $|u|<\frac12$.

(b) \emph{Step 1: the factor $\mathcal J$.} Write $c=\kappa_s$,
$c'=c+\kappa$ and $\Sigma=\Sigma_s$. Since
$c'-\beta s=(\kappa-\beta)s+11\kappa$, the choice of $\kappa$ gives
$\rho-1=\kappa+\nu_s$ with $\nu_s=\frac{22\kappa}{s(1+2\beta)}$. As $u$ and
$v$ have the sign of $d$, $u(u+v)\ge u^2$. So
$\mathcal J\le1-(\kappa+\nu_s)W'$ with $z'=4u^2$ and $W'=\frac{z'}{1-z'}$.
Also $\mathcal J>0$ by (a), since $0<r<1$ and $\bar r$, $1-\bar r$ are not
both 0. So $\log\mathcal J\le-(\kappa+\nu_s)W'$.

\emph{Step 2: the ratio $\Phi^{c'}_s/\Phi^c_s$.} The logarithm
$\log(\Phi^{c'}_s(b)/\Phi^c_s(b))$ is the integral over $\tau\in[c,c']$ of
the derivative in Lemma~\ref{lem:A-cderiv}. Its bound increases in $z$ and
decreases in $\Sigma$, while $z$ decreases and $\Sigma$ increases in $\tau$.
So we may take the bound at $\tau=c$, where $X+Y=\Sigma$, and multiply it by
the length $\kappa$ of the interval. Put
$W=\frac z{1-z}$. Then $\frac1{1-z}=1+W$, $\frac{1+z}{1-z}=1+2W$ and
$-\log(1-z)\le W$, so
\begin{equation}\label{eq:A-logD}
 \log \Delta_s(b)\le\kappa\Bigl[W+\frac{1+2W}\Sigma+\frac{2(1+2W)(1+W)}{3\Sigma^2}\Bigr]
 -(\kappa+\nu_s)W'.
\end{equation}
\emph{Step 3: 3 facts.} (i) $X,Y\ge1+c$, since $b\ge1$ and $s-b\ge1$. So
$XY\ge(1+c)\Sigma/2$ and $1+W=\frac{\Sigma^2}{4XY}\le\frac\Sigma{2(1+c)}$.
(ii) $z'=\lambda z$ with $\lambda=(\frac\Sigma{\Sigma+2\kappa})^2
\ge1-\frac{4\kappa}\Sigma$. So $W-W'=\frac{(1-\lambda)z}{(1-z)(1-\lambda z)}
\le(1-\lambda)W(1+W)\le\frac{2\kappa}{1+c}W$ by (i). As $c\ge11\kappa$, this
gives $W'\ge\frac9{11}W$. (iii) $\Sigma\ge s$ and $1+c>\kappa s$. If $W>0$
then $s\ne2b$, and with $b\ge1$, $s-b\ge1$ this forces $s\ge3$, so
$\Sigma^2\ge3s$.

\emph{Step 4: collect the terms.} By (ii) and (iii),
$\kappa(W-W')\le\frac{2\kappa^2}{1+c}W\le\frac{2\kappa}sW$
and $\frac{2\kappa W}\Sigma\le\frac{2\kappa}sW$. Since
$(1+2W)(1+W)=1+3W+2W^2$, $\kappa$ times the last term in the bracket of
\eqref{eq:A-logD} is
$\frac{2\kappa}{3\Sigma^2}+\frac{2\kappa W}{\Sigma^2}
+\frac{4\kappa W^2}{3\Sigma^2}$. Here
$\frac{2\kappa W}{\Sigma^2}\le\frac{2\kappa}{3s}W$ by (iii), and
$\frac{4\kappa W^2}{3\Sigma^2}\le\frac{4\kappa W(1+W)}{3\Sigma^2}
\le\frac{2\kappa W}{3\Sigma(1+c)}\le\frac{2\kappa}{3s}W$ by (i). With
$-\nu_sW'\le-\frac9{11}\nu_sW$ this gives
\[
 \log \Delta_s(b)\le\frac\kappa\Sigma+\frac{2\kappa}{3\Sigma^2}
 +W\Bigl[\frac{16\kappa}{3s}-\frac{18\kappa}{s(1+2\beta)}\Bigr],
\]
and the bracket is negative because $1+2\beta<2$ for $\varepsilon<\frac14$,
so $18/(1+2\beta)>9>16/3$.

(c) Apply Lemma~\ref{lem:A-compare}(a) with $s_0=t+1$, $F_s=\Phi^{\kappa_s}_s$
and $K_s=\exp(\frac\kappa s+\frac{2\kappa}{3s^2})$. This is allowed, since
$F_N(h)=1$, $F_s>0$ on $\mathcal S_s$, and $\mathcal L_sF_{s+1}=\Delta_sF_s\le
K_sF_s$ on $\mathcal S_s$ by (b) and $\Sigma_s\ge s$. Since $1\in\mathcal S_{t+1}$ and
$\sum_{s>t}s^{-2}\le\frac1t$, we get
$V_{t+1}(1)\le e^{\kappa[H_{N-1}-H_t+2/(3t)]}\,\Phi^{\kappa_{t+1}}_{t+1}(1)$.
Now $\log(\Phi^{\kappa_{t+1}}_{t+1}(1)/\pi_t)$ is the integral over
$\tau\in[0,\kappa_{t+1}]$ of the derivative in Lemma~\ref{lem:A-cderiv}, whose
bound is largest at $\tau=0$. There $X=1$, $Y=t$ and
$\frac1{1-z}=\frac{(t+1)^2}{4t}$, and the bound is
$\log\frac{(t+1)^2}{4t}+\frac{t^2+1}{2t(t+1)}+\frac{t^2+1}{12t^2}
\le\log\frac{(t+1)^2}{4t}+\frac23$.

(d) Insert (c) into Proposition~\ref{prop:first-mutation}(a), and use
$\kappa_{t+1}=\kappa(t+11)$ and $\pi_t=c_1w_t$.
\end{proof}

\begin{proof}[Proof of Theorem~\ref{thm:erw-bounds}]
(a) By Proposition~\ref{prop:first-mutation}(a) and
Lemma~\ref{lem:A-polya}(d),
$C_N\ge\varepsilon c_1\sum_tw_t(1-\varepsilon)^{t-1}$. Now
$(1-\varepsilon)^{t-1}\ge1-(t-1)\varepsilon$ and
$\sum_tw_t(t-1)=2-\frac6{h+2}$ by Proposition~\ref{prop:first-mutation}(c).
So $C_N\ge\varepsilon c_1(1-2\varepsilon)$, and $c_1=\frac4{N+2}$.

(b) \emph{Step 1: the constants.} Let $N\ge64$ and
$\varepsilon(1+\log N)\le\frac1{100}$. Then
$\varepsilon<0.002$ and $\kappa\le2\varepsilon(1+8\varepsilon)\le2.1\varepsilon$.
Since $\log\frac{(t+1)^2}{4t}\ge0$ and $\frac23(t+11)>1+\log t\ge H_t$, we
have $\mathcal G(t)>0$. Since $\frac{(t+1)^2}{4t}\le\frac{t+1}2$,
$\frac23<\log2$ and $H_t\ge\frac2{3t}$, also
$\mathcal G(t)\le(t+11)\log(t+1)\le12t\log(t+1)$. So
$\kappa\mathcal G(t)\le25.2\,\varepsilon t\log(t+1)$, and this is at most
$\frac t2\log2$ for $t\le h$, since then
$25.2\,\varepsilon\log(t+1)\le25.2\,\varepsilon\log N\le0.252<\frac12\log2$.

\emph{Step 2: the sum over $t$.} Since $e^y-1\le ye^y$ for $y\ge0$, Step~1
gives $e^{\kappa\mathcal G(t)}-1\le\kappa\mathcal G(t)2^{t/2}$. By
Proposition~\ref{prop:A-super}(d), $(1-\varepsilon)^{t-1}\le1$,
$\sum_tw_t=1$ and $w_t\le2t2^{-t}$ (Proposition~\ref{prop:first-mutation}(c)),
\[
 C_N\le\varepsilon c_1e^{\kappa H_{N-1}}\Bigl(1+50.4\,\varepsilon
 \sum_{t\ge1}t^2\log(t+1)2^{-t/2}\Bigr),
\]
and we bound the series as follows. The ratio of its consecutive terms is
$(1+\frac1t)^2\frac{\log(t+2)}{\log(t+1)}2^{-1/2}$, which decreases in $t$ and
is below $\frac34$ at $t=60$. So the terms with $t>60$ add up to less than 4
times the term at $t=61$, which is below $10^{-4}$. The first 60 terms add up
to $102.5864$ to 4 decimals (computed in 30-digit arithmetic by
\texttt{verify\_elephant\_crossing.py}), so the series is below $102.59$.

\emph{Step 3: the factor $e^{\kappa H_{N-1}}$.} Finally
$H_{N-1}\le1+\log N$, so
$\kappa H_{N-1}\le2\varepsilon\log N+2\varepsilon+16\varepsilon^2(1+\log N)
\le2\varepsilon\log N+3\varepsilon$. With $1+y\le e^y$ this proves (b),
since $3+50.4\cdot102.59<5200$.
\end{proof}
